\scanpage{48}
Montrer que les deux images inverses de $H'$ dans $G''=G\times_S\unsure{S'|S''}$ (où $S''=S'\times_S S'$) sont les $\overline{\mathrm{m}}$s.

\deleted{On peut supposer $S$ noethérien et $=$ affine \inserted{d'ann. $A$}, et $S'$ affine sur $S$ [en remplaçant $S'$ par \deleted{\unsure{la}} somme de ses ouverts d'un rec. affine] ; d'ann. $A'$. Or $A'=\lim A'_i$, les $A'_i$ étant des sous-algèbres de \deleted{\unsure{$A_i$}} de type fini de $A'$, on $\unsure{A'_i|A'}\otimes_A A'=\varinjlim A'_i\otimes_A A'$. Comme $H'$ provient d'un $H'_i$ défini sur un $A'_i$, également rigide [?],}
\margin{\deleted{O.K. dans le cas de schémas abéliens ou multiplicatifs ou \deleted{\unsure{anti}} \unsure{anadditifs|antiadditifs|additifs}}}
\deleted{il suffit donc de prouver que les deux images inverses de $H'_i$ en $G''_i=G\times_A A''_i$\deleted{$\unsure{i}$}, où $A''_i=A'_i\otimes A'_i$, sont les $\overline{\mathrm{m}}$s. \deleted{\unsure{Or}} Comme $A''_i$ est noethérien, il suffit pour cela de :}

Or elles coïncident sur la diagonale de $S'\times_S S'$, donc sur \deleted{toute} composante connexe de $S'\times_S S'$ qui contient la dite par le th. Comme les fibres de $S''=S'\times_S S'$ sont connexes [car celles de $S'$ sont \unsure{abst|géom.} connexes] tt compos conn. de $S''$ contient au moins une fibre, donc
% archive sheet number: 47
